\documentclass[14pt, a4paper]{extarticle}
\usepackage[T2A]{fontenc}
\usepackage[utf8]{inputenc}
\usepackage[english, russian]{babel}
\usepackage{tempora}
\usepackage{longtable}
\usepackage{array}
\usepackage{ltablex}
\usepackage{ragged2e}
\usepackage{graphicx}
\usepackage{calc}
\usepackage[left=30mm, right=15mm, top=20mm, bottom=20mm, footskip=12mm]{geometry}

\setlength{\parindent}{0pt}
\setlength{\parskip}{6pt plus 2pt minus 1pt}

\newcolumntype{P}[1]{>{\RaggedRight\arraybackslash}p{#1 - 2\tabcolsep - 2\arrayrulewidth}}
\newcolumntype{C}[1]{>{\centering\arraybackslash}p{#1 - 2\tabcolsep - 2\arrayrulewidth}}


\begin{document}

\pagestyle{empty}

% ------------------------- Заголовок документа -------------------------

\begin{center}
    {\small МИНИСТЕРСТВО НАУКИ И ВЫСШЕГО ОБРАЗОВАНИЯ РОССИЙСКОЙ ФЕДЕРАЦИИ}

    \resizebox{\textwidth}{!}{\fontsize{8pt}{10pt}\selectfont ФЕДЕРАЛЬНОЕ
        ГОСУДАРСТВЕННОЕ АВТОНОМНОЕ ОБРАЗОВАТЕЛЬНОЕ УЧРЕЖДЕНИЕ ВЫСШЕГО ОБРАЗОВАНИЯ}

    {\small \textbf{«МОСКОВСКИЙ ПОЛИТЕХНИЧЕСКИЙ УНИВЕРСИТЕТ»}} \\
    {\small \textbf{(МОСКОВСКИЙ ПОЛИТЕХ)}}
\end{center}

\begin{tabularx}{\textwidth}{@{} p{0.53\textwidth} p{0.43\textwidth} @{}}
     &
    \raggedright
    УТВЕРЖДАЮ                          \\
    заведующая кафедрой                \\
    «Инфокогнитивные технологии»       \\[0.3cm]

    \rule{3cm}{0.4pt} / Е.~А.~Пухова / \\[0.1cm]

    «\rule{0.8cm}{0.4pt}» \rule{3cm}{0.4pt} 2026~г.
\end{tabularx}

% -------------------------- Общая информация --------------------------

\begin{center}

    { \large \textbf{ЗАДАНИЕ}}

    \textbf{на выполнение курсовой работы (проекта)}

    Мижитдоржиеву Гэлэгу Намдаковичу, % нужно указать в дательном падеже (задание кому?)

\end{center}

обучающемуся группы 241-327,

направления подготовки 09.03.01 «Информатика и вычислительная техника»

по дисциплине «Разработка веб-приложений»

на тему «Разработка веб-приложения для управления и мониторинга CI/CD-пайплайнов»

1. Исходные данные к работе (проекту): информационные ресурсы в сети
интернет, научные публикации в открытой печати.

2. Содержание задания по курсовой работе (проекту) -- перечень вопросов,
подлежащих разработке:

% ------------------------------- Задачи -------------------------------

{
\renewcommand{\arraystretch}{1.1}
\footnotesize
\begin{longtable}{|P{0.50\textwidth}|C{0.11\textwidth}|C{0.15\textwidth}|C{0.24\textwidth}|}

    \hline
    \textbf{Разрабатываемый вопрос}                                                                  & \textbf{Объем, \%} & \textbf{Срок} & \textbf{Примечание} \\
    \hline
    \endhead

    \hline
    \endfoot

    \multicolumn{4}{|P{\textwidth}|}{\textbf{Раздел 1. Анализ предметной области}}                                                                              \\ \hline
    Задача 1.1. Обзор существующих CI/CD-систем и инструментов автоматизации                               & 10                 & 22.03.2026    &                     \\ \hline
    Задача 1.2. Анализ программных средств разработки и обоснование технологического стека                 & 7                  & 26.03.2026    &                     \\ \hline
    Задача 1.3. Формулировка целей, задач и технических требований к веб-приложению                        & 8                  & 30.03.2026    &                     \\ \hline

    \multicolumn{4}{|P{\textwidth}|}{\textbf{Раздел 2. Проектирование веб-приложения}}                                                                          \\ \hline
    Задача 2.1. Анализ целевой аудитории и вариантов использования (Use Case, User Story)                  & 6                  & 04.04.2026    &                     \\ \hline
    Задача 2.2. Проектирование модели данных (ER-диаграмма, физическая схема PostgreSQL)                  & 7                  & 09.04.2026    &                     \\ \hline
    Задача 2.3. Проектирование архитектуры безопасности (OAuth2, JWT, шифрование Fernet)                   & 6                  & 14.04.2026    &                     \\ \hline
    Задача 2.4. Разработка макетов пользовательского интерфейса (Wireframe)                                & 6                  & 18.04.2026    &                     \\ \hline

    \multicolumn{4}{|P{\textwidth}|}{\textbf{Раздел 3. Разработка веб-приложения}}                                                                              \\ \hline
    Задача 3.1. Разработка асинхронного REST API на FastAPI и ORM-моделей БД                      & 8                  & 25.04.2026    &                     \\ \hline
    Задача 3.2. Реализация модуля аутентификации (JWT) и подсистемы безопасного хранения токенов провайдеров & 7                  & 01.05.2026    &                     \\ \hline
    Задача 3.3. Разработка сервиса интеграции с GitHub Actions API и синхронизации пайплайнов              & 8                  & 07.05.2026    &                     \\ \hline
    Задача 3.4. Реализация диспетчеризации сборок (Re-run) и аналитического модуля расчета KPI             & 9                  & 13.05.2026    &                     \\ \hline
    Задача 3.5. Разработка клиентского интерфейса на React (SPA) и интеграция с бэкендом                   & 8                  & 19.05.2026    &                     \\ \hline

    \multicolumn{4}{|P{\textwidth}|}{\textbf{Раздел 4. Оформление итогов работы}}                                                                               \\ \hline
    Задача 4.1. Контейнеризация системы (Docker Compose, Nginx) и публикация репозитория                  & 4                  & 22.05.2026    &                     \\ \hline
    Задача 4.2. Развертывание сервисов на облачной платформе Render                               & 3                  & 24.05.2026    &                     \\ \hline
    Задача 4.3. Оформление расчетно-пояснительной записки                                        & 3                  & 26.05.2026    &                     \\ \hline
\end{longtable}
}

% -------------- Сведения о выдаче и принятии задания -------------------

Руководитель курсовой работы (проекта): \\ старший преподаватель кафедры
«Инфокогнитивные технологии»

{
\renewcommand{\arraystretch}{2}
\noindent
\begin{tabularx}{\textwidth}{@{} P{0.45\textwidth} P{0.3\textwidth} P{0.25\textwidth} @{}}

    «30»~марта 2026~г. &
    \centering \rule{3.5cm}{0.4pt} \par \footnotesize (подпись) &
    \hfill А.~С.~Кружалов                                         \\
\end{tabularx}
\begin{tabularx}{\textwidth}{@{} P{0.45\textwidth} P{0.55\textwidth} @{}}
    Дата выдачи задания           &
    \hfill «30»~марта 2026~г. \\

    Дата сдачи выполненной работы &
    \hfill «22»~сентября 2026~г. \\
\end{tabularx}
}

Задание принял к исполнению

{
\renewcommand{\arraystretch}{1.5}
\begin{tabularx}{\textwidth}{@{} P{0.45\textwidth} P{0.3\textwidth} P{0.25\textwidth} @{}}

    «30»~марта 2026~г. &
    \centering \rule{3.5cm}{0.4pt} \par \footnotesize (подпись) &
    \hfill Г.~Н.~Мижитдоржиев                                            \\
\end{tabularx}
}

\end{document}
